% !TEX root = hw11-answers-graded.tex
\chead{\textbf{Exercises week 11}}
\begin{document}
\flushleft\textbf{Topic: Examples of linear SCMs and acyclifications}

\begin{enumerate}[label=11.\arabic*]
	\item %($3$ points)
	We will prove that the SCM $M$ of the equilibrium behaviour of a damped coupled harmonic oscillator in Example \ref*{ex:HarmonicOscillator} is simple:

	\begin{enumerate}
		% \item Prove that for any\footnote{It's more straightforward to see this in the general case.} SCM satisfying the assumptions of Proposition \ref*{prp:linear_scms} that invertibility of $1_{V}-B_{VV}$ ensures that the SCM has a unique solution.
		\item Show that $M$ is essentially linear (Definition \ref*{def:linear_scms}).
		\item \label{eq:tridiag} Show that the determinant of the matrix $1_{V}-B_{VV}$ (where $1_V$ is the $|V|\times |V|$ identity matrix) is non-zero, and hence, that the matrix is invertible.

		\emph{Hint: you can use that the determinant of a tridiagonal matrix of the following form is given by the expression on the r.h.s.
		% \footnote{Bonus points for who comes up with an elegant proof of this identity that does not rely on solving the recurrence relationship for the determinant of tridiagonal matrices.}
		:
		$$\det \begin{pmatrix}
				k_0+k_1 & -k_1      &           &          &               \\
				-k_1    & k_1 + k_2 & -k_2      &          &               \\
				        & -k_2      & k_2 + k_3 & \ddots   &               \\
				        &           & \ddots    & \ddots   & -k_{d-1}      \\
				        &           &           & -k_{d-1} & k_{d-1} + k_d \\
			\end{pmatrix} = \sum_{i=0}^d \prod_{j=0 \atop j\ne i}^d k_j
		$$
		In order to obtain the above form, factorize the matrix $1_{V}-B_{VV}=DT$, where $T$ is of the above form and $D$ is diagonal. Use $\det DT=\det D \det T$.}
		\item Show that for for any $T \subseteq V$, the matrix $1_{V \setminus T}-B_{(V \setminus T)(V \setminus T)}$ is invertible, and conclude that the SCM $M$ is simple.\\
		\emph{Hint: You can use the fact that the determinant of a block-diagonal matrix is the product of the determinant of the constituent matrices.}
	\end{enumerate}
	\begin{ungradedsolution}
		\begin{enumerate}
			\item We have
			$$f(q, k, \ell) = B_{VV} \cdot q + c_V$$
			where
			$$B_{VV} := \mat{0 & \tfrac{k_1}{k_0+k_1}& 0 & 0 & \hdots & 0 \\ \tfrac{k_1}{k_1+k_2} & 0 & \tfrac{k_2}{k_1+k_2} & 0 & &\vdots \\ 0 & \tfrac{k_2}{k_2+k_3} & 0 &  \tfrac{k_3}{k_2+k_3} & \ddots & \vdots\\ \vdots& & & & & \vdots \\ \vdots& & & \ddots & 0 &  \tfrac{k_{d-1}}{k_{d-2}+k_{d-1}} \\ 0 & \hdots &\hdots & \hdots & \tfrac{k_{d-1}}{k_{d-1}+k_d} & 0 \\}$$
			and
			$$c_V := \mat{\tfrac{k_0q_0+k_0\ell_0-k_1\ell_1}{k_0+k_1} \\ \tfrac{k_1\ell_1-k_2\ell_2}{k_1+k_2} \\ \vdots \\ \tfrac{k_dq_{d+1}+k_{d-1}\ell_{d-1}-k_d\ell_d}{k_{d-1}+k_d}}.$$
			\item
			\begin{align*}
				 & 1_V-B_{VV}                                                                       \\
				 & =\begin{pmatrix}
					    1                & -k_1/(k_1+k_0)   & 0                & 0                & ... \\
					    -k_1 / (k_2+k_1) & 1                & -k_2 / (k_2+k_1) & 0                & ... \\
					    0                & -k_2 / (k_3+k_2) & 1                & -k_3 / (k_3+k_2) & ... \\
					    ...              & ...              & ...              & ...              & ...
				    \end{pmatrix} \\
				 & =\begin{pmatrix}
					    1/(k_1+k_0) & 0            & 0             & 0   & ... \\
					    0           & 1/ (k_2+k_1) &               & 0   & ... \\
					    0           & 0            & 1 / (k_3+k_2) & 0   & ... \\
					    ...         & ...          & ...           & ... & ...
				    \end{pmatrix}
				\begin{pmatrix}
					k_0+k_1 & -k_1      & 0         & \dots \\
					-k_1    & k_1 + k_2 & -k_2      & \dots \\
					0       & -k_2      & k_2 + k_3 & \dots \\
					\dots   & \dots     & \dots     & \dots
				\end{pmatrix}
			\end{align*}
			The determinant of the first matrix is the product of the diagonals, the second is the given by the hint, so
			\begin{align*}
				\det (1_V - B_{VV}) = \left(\prod_{i=1}^d \frac{1}{k_{i-1}+k_i}\right)\left(\sum_{i=0}\prod_{\substack{j = 0 \\j\neq i}}^d k_j\right)
			\end{align*}
			which is clearly positive for positive $k_i$.

			\item For any $T \subseteq V$, the matrix $(1_{V}-B_{VV})_{(V \setminus T)(V \setminus T)} = 1_{V \setminus T}-B_{(V \setminus T)(V \setminus T)}$ is a block-diaginal matrix, where each block is either a diagonal or a tridiagonal matrix. (Check this!) The determinant of a block-diagonal matrix equals the product of the determinants of the constituent blocks. Using the same analysis as in the previous exercise, it follows that $\det(1_{V \setminus T}-B_{(V \setminus T)(V \setminus T)}) > 0$, hence $1_{V \setminus T}-B_{(V \setminus T)(V \setminus T)}$ is invertible. From Proposition \ref*{prp:linear_scms} we know that for every $T\subseteq V$ the SCM $M_{\doit (T)}$ is essentially uniquely solvable, and from the definition of simple SCMs (Definition \ref*{def:simple_scm}) we conclude that $M$ is simple.
		\end{enumerate}
	\end{ungradedsolution}

	\item %(1 point) 
	Consider the dynamical system and the SCM given in Example \ref*{ex:PriceSupplyDemandDynamical} (Price, supply and demand) of the lecture notes.

	Show for which of the 8 subsets $T \subseteq \{D, S, R\}$, the intervened SCM $M_{\text{do}(T)}$ is essentially uniquely solvable. You can use the condition for unique solvability of linear SCMs from Proposition \ref*{prp:linear_scms} of the lecture notes.
	% \item For all possible interventions $T \subseteq \{D, S, R\}$, show whether the dynamical system has an equilibrium and, if it exists, compute this equilibrium.
	% We see that when the intervened SCM is essentially uniquely solvable, its solution indeed equals the equilibrium of the intervened dynamical system.
	\begin{ungradedsolution}
		We have a linear system:
		\begin{align*}
			\mat{f_D \\f_S\\f_R} = \mat{0 & 0 & \beta_D \\ 0 & 0 & \beta_S \\ 1 & -1 & 1}\cdot \mat{d\\s\\r}+\mat{e_D\\e_S\\0}
		\end{align*}
		so by Proposition \ref*{prp:linear_scms} there exists a unique solution function for $M_{\doit(T)}$ iff
		\begin{multline*}
			\det\left(\left(\mat{1&0&0\\0&1&0\\0&0&1} - \mat{0 & 0 & \beta_D \\ 0 & 0 & \beta_S \\ 1 & -1 & 1}\right)_{(V\setminus T)(V\setminus T)}\right) \\
			= \det\left(\mat{1 & 0 & -\beta_D \\ 0 & 1 & -\beta_S \\ -1 & 1 & 0}_{(V\setminus T)(V\setminus T)}\right)
		\end{multline*}
		is positive. For the following $T$ we have the determinant: $\{D\}: \beta_S$, $\{S\}:-\beta_D$, $\{R\}: 1$, $\{D, S\}: 0$, $\{D, R\}: 1$, $\{S, R\}: 1$, $\emptyset: \beta_s-\beta_D$. By assumption, $\beta_D < 0$ and $\beta_S> 0$, so $\beta_S-\beta_D > 0$. Note that for $T=\{D, S, R\}$ the structural equations are already a (unique) solution function, as they don't depend on $(d, s, r)$. Hence, for all subsets $T\subseteq\{D, S, R\}$ except for $\{D, S\}$, the SCM $M_{\doit(T)}$ is essentially uniquely solvable.
	\end{ungradedsolution}

	\item %($5$ points) 
	Given the following SCM:
	$$M = \left( J=\emptyset, V=\{A, B, C, D, E, F\}, W=\{A', B', C', D', E', F'\}, \R^{12}, P, f \right),$$
	where $P$ consists of the product of six independent normal distributions with means and standard deviations $\mu_A, \sigma_A, \mu_B, ..., \sigma_F$, and
	\[
		\begin{pmatrix}
			f_A(x_V, \varepsilon_W) \\ f_B(x_V, \varepsilon_W) \\ f_C(x_V, \varepsilon_W) \\ f_D(x_V, \varepsilon_W) \\ f_E(x_V, \varepsilon_W) \\ f_F(x_V, \varepsilon_W)
		\end{pmatrix}
		=
		\begin{pmatrix}
			0           & 0           & \alpha_{ca} & \alpha_{da} & 0           & 0 \\
			\alpha_{ab} & 0           & 0           & 0           & \alpha_{eb} & 0 \\
			0           & \alpha_{bc} & 0           & 0           & 0           & 0 \\
			0           & 0           & 0           & 0           & 0           & 0 \\
			0           & 0           & 0           & 0           & 0           & 0 \\
			0           & 0           & \alpha_{cf} & 0           & 0           & 0 \\
		\end{pmatrix}
		\begin{pmatrix}
			x_A \\ x_B \\ x_C \\ x_D \\ x_E \\ x_F
		\end{pmatrix}
		+
		\begin{pmatrix}
			\varepsilon_A \\ \varepsilon_B \\ \varepsilon_C \\ \varepsilon_D \\ \varepsilon_E \\ \varepsilon_F
		\end{pmatrix}
	\]
	for nonzero real numbers $\alpha_.$ such that $\alpha_{ab}\alpha_{bc}\alpha_{ca} \neq 1$. Note that we have a slightly different notation as in the lecture notes, as we refer to $x_V$ as $x$, and refer to $x_W$ as $\varepsilon$.
	\begin{enumerate}
		\item Construct the acyclification $M^{\text{acy}}$.\\
		\emph{Hint: use Proposition \ref*{prp:linear_scms}.}
		\item Give the graphs $G^+(M)$ and $G^+(M^{\text{acy}})$.
		Give also the causal graphs $G(M)$ and $G(M^{\text{acy}})$.
		\item Show that $M^{\text{acy}}$ is observationally equivalent to $M$: calculate the solution functions $g(\varepsilon_W)$ and $\tilde{g}(\varepsilon_W)$ of $M$ and $M^{\text{acy}}$ respectively, and ague why it follows that  distributions $P_M(X_V) = P_{M^{\text{acy}}}(X_V)$. For this last step you may refer to a result from the lecture notes, but you may not use Proposition \ref*{prp:scm_acyclification_properties}.
		\item Show that $M^{\text{acy}}$ is not interventionally equivalent to $M$ by giving an example of a set $T\subseteq V$ such that $P_M(X_{V\setminus T} \given \doit(X_T=x_t)) \neq P_{M^\text{acy}}(X_{V\setminus T} \given \doit(X_T=x_t))$ (prove this inequality).
		\item Show the following equivalences for this particular example, in the sense that both statements are either true or false:
		\begin{enumerate}
			\item $D \perp_{G(M)}^\sigma F \mid B \iff D \perp_{G(M^{\text{acy}})}^d F \mid B$
			\item $D \perp_{G(M)}^\sigma E \mid A \iff D \perp_{G(M^{\text{acy}})}^d E \mid A$
			\item $D \perp_{G(M)}^\sigma E \iff D \perp_{G(M^{\text{acy}})}^d E$
		\end{enumerate}
		Do this by giving $\sigma$/$d$-active paths or walks, or prove the $\sigma$/$d$-separation.
	\end{enumerate}

	\begin{gradedsolution}
		\begin{enumerate}
			\item \points{3} We have strongly connected component $\{A, B, C\}$. Applying Proposition \ref*{prp:linear_scms} with $L = \{A, B, C\}$,
			\begin{equation*}
				B_{VV}(x_J, \varepsilon) = \begin{pmatrix}
					0           & 0           & \alpha_{ca} & \alpha_{da} & 0           & 0 \\
					\alpha_{ab} & 0           & 0           & 0           & \alpha_{eb} & 0 \\
					0           & \alpha_{bc} & 0           & 0           & 0           & 0 \\
					0           & 0           & 0           & 0           & 0           & 0 \\
					0           & 0           & 0           & 0           & 0           & 0 \\
					0           & 0           & \alpha_{cf} & 0           & 0           & 0 \\
				\end{pmatrix}
				\text{ and }
				c_V(x_J, \varepsilon) = \begin{pmatrix}
					\varepsilon_A \\ \varepsilon_B \\ \varepsilon_C \\ \varepsilon_D \\ \varepsilon_E \\ \varepsilon_F
				\end{pmatrix},
			\end{equation*}
			the solution function $g_L$ is given by
			\begin{align*}
				g_L(x_{\{D, E, F\}}, \varepsilon) = \left(\mat{1 & 0           & 0              \\0&1&0\\0&0&1} - \mat{0 & 0 & \alpha_{ca}  \\
				\alpha_{ab}                                      & 0           & 0              \\
				0                                                & \alpha_{bc} & 0}\right)^{-1} \\
				\cdot \left(\mat{\alpha_{da}                     & 0           & 0              \\0 & \alpha_{eb} & 0\\ 0 & 0 & 0}\mat{x_D\\ x_E\\ x_F} + \mat{\varepsilon_A\\\varepsilon_B\\\varepsilon_C}\right)
			\end{align*}
			Computing the inverse, we have
			\begin{equation*}
				\left(\mat{1 & 0           & 0              \\0&1&0\\0&0&1} - \mat{0 & 0 & \alpha_{ca}  \\
					\alpha_{ab}                                      & 0           & 0              \\
					0                                                & \alpha_{bc} & 0}\right)^{-1} = \frac{1}{1-\alpha_{ca}\alpha_{ab}\alpha_{bc}}\begin{pmatrix}
					1                      & \alpha_{bc}\alpha_{ca} & \alpha_{ca}            \\
					\alpha_{ab}            & 1                      & \alpha_{ab}\alpha_{ca} \\
					\alpha_{ab}\alpha_{bc} & \alpha_{bc}            & 1
				\end{pmatrix}
			\end{equation*}
			so
			\begin{align*}
				g_L(x_{\{D, E, F\}}, \varepsilon) = & \frac{1}{1-\alpha_{ca}\alpha_{ab}\alpha_{bc}}
				\begin{pmatrix}
					\alpha_{da}                       & \alpha_{bc}\alpha_{ca}\alpha_{eb} & 0 \\
					\alpha_{ab}\alpha_{da}            & \alpha_{eb}                       & 0 \\
					\alpha_{ab}\alpha_{bc}\alpha_{da} & \alpha_{bc}\alpha_{eb}            & 0
				\end{pmatrix}
				\begin{pmatrix}x_D \\ x_E\\ x_F \end{pmatrix}+                                      \\
				                                    & \frac{1}{1-\alpha_{ca}\alpha_{ab}\alpha_{bc}}
				\begin{pmatrix}
					1                      & \alpha_{bc}\alpha_{ca} & \alpha_{ca}            \\
					\alpha_{ab}            & 1                      & \alpha_{ab}\alpha_{ca} \\
					\alpha_{ab}\alpha_{bc} & \alpha_{bc}            & 1
				\end{pmatrix}
				\begin{pmatrix}
					\varepsilon_A \\
					\varepsilon_B \\
					\varepsilon_C
				\end{pmatrix}.
			\end{align*}
			In the SCM, we replace the equations $f_{\{A,B,C\}}$ with this solution.
			\item \points{2} $G^+(M)$:
			\begin{center}
				\begin{tikzpicture}[scale=1, transform shape]
					\node[ndout] (A) at (0, 0) {$A$};
					\node[ndout] (B) at (2, 0) {$B$};
					\node[ndout] (C) at (4, 0) {$C$};
					\node[ndout] (D) at (0, -2) {$D$};
					\node[ndout] (E) at (2, -2) {$E$};
					\node[ndout] (F) at (4, -2) {$F$};
					\node[ndlat] (epsA) at (0, 2) {$A'$};
					\node[ndlat] (epsB) at (2, 2) {$B'$};
					\node[ndlat] (epsC) at (4, 2) {$C'$};
					\node[ndlat] (epsD) at (0, -4) {$D'$};
					\node[ndlat] (epsE) at (2, -4) {$E'$};
					\node[ndlat] (epsF) at (4, -4) {$F'$};
					\draw[arlout] (epsA) to (A);
					\draw[arlout] (epsB) to (B);
					\draw[arlout] (epsC) to (C);
					\draw[arlout] (epsD) to (D);
					\draw[arlout] (epsE) to (E);
					\draw[arlout] (epsF) to (F);
					\draw[arout] (D) to (A);
					\draw[arout] (E) to (B);
					\draw[arout] (A) to (B);
					\draw[arout] (B) to (C);
					\draw[arout, bend right] (C) to (A);
					\draw[arout] (C) to (F);
				\end{tikzpicture}
			\end{center}
			$G^+(M^{\text{acy}})$:
			\begin{center}
				\begin{tikzpicture}[scale=1, transform shape]
					\node[ndout] (A) at (0, 0) {$A$};
					\node[ndout] (B) at (2, 0) {$B$};
					\node[ndout] (C) at (4, 0) {$C$};
					\node[ndout] (D) at (0, -2) {$D$};
					\node[ndout] (E) at (2, -2) {$E$};
					\node[ndout] (F) at (4, -2) {$F$};
					\node[ndlat] (epsA) at (0, 2) {$A'$};
					\node[ndlat] (epsB) at (2, 2) {$B'$};
					\node[ndlat] (epsC) at (4, 2) {$C'$};
					\node[ndlat] (epsD) at (0, -4) {$D'$};
					\node[ndlat] (epsE) at (2, -4) {$E'$};
					\node[ndlat] (epsF) at (4, -4) {$F'$};
					\draw[arlout] (epsA) to (A);
					\draw[arlout] (epsA) to (B);
					\draw[arlout] (epsA) to (C);
					\draw[arlout] (epsB) to (A);
					\draw[arlout] (epsB) to (B);
					\draw[arlout] (epsB) to (C);
					\draw[arlout] (epsC) to (A);
					\draw[arlout] (epsC) to (B);
					\draw[arlout] (epsC) to (C);
					\draw[arlout] (epsD) to (D);
					\draw[arlout] (epsE) to (E);
					\draw[arlout] (epsF) to (F);
					\draw[arout] (D) to (A);
					\draw[arout] (D) to (B);
					\draw[arout] (D) to (C);
					\draw[arout] (E) to (B);
					\draw[arout] (E) to (A);
					\draw[arout] (E) to (C);
					\draw[arout] (C) to (F);
				\end{tikzpicture}
			\end{center}
			$G(M)$:
			\begin{center}
				\begin{tikzpicture}[scale=1, transform shape]
					\node[ndout] (A) at (0, 0) {$A$};
					\node[ndout] (B) at (2, 0) {$B$};
					\node[ndout] (C) at (4, 0) {$C$};
					\node[ndout] (D) at (0, -2) {$D$};
					\node[ndout] (E) at (2, -2) {$E$};
					\node[ndout] (F) at (4, -2) {$F$};
					\draw[arout] (D) to (A);
					\draw[arout] (E) to (B);
					\draw[arout] (A) to (B);
					\draw[arout] (B) to (C);
					\draw[arout, bend right] (C) to (A);
					\draw[arout] (C) to (F);
				\end{tikzpicture}
			\end{center}
			$G(M^{\text{acy}})$:
			\begin{center}
				\begin{tikzpicture}[scale=1, transform shape]
					\node[ndout] (A) at (0, 0) {$A$};
					\node[ndout] (B) at (2, 0) {$B$};
					\node[ndout] (C) at (4, 0) {$C$};
					\node[ndout] (D) at (0, -2) {$D$};
					\node[ndout] (E) at (2, -2) {$E$};
					\node[ndout] (F) at (4, -2) {$F$};
					\draw[arlat] (A) edge (B);
					\draw[arlat] (B) edge (C);
					\draw[arlat,bend left] (A) edge (C);
					\draw[arout] (D) to (A);
					\draw[arout] (D) to (B);
					\draw[arout] (D) to (C);
					\draw[arout] (E) to (B);
					\draw[arout] (E) to (A);
					\draw[arout] (E) to (C);
					\draw[arout] (C) to (F);
				\end{tikzpicture}
			\end{center}
			\item \points{2} 1 point for computing $g$, and 1 point for computing $\tilde{g}$. Writing $\Delta := 1-\alpha_{ca}\alpha_{ab}\alpha_{bc}$ we have
			\begin{equation*}
				g(\varepsilon) =
				\frac{1}{\Delta}
				\begin{pmatrix}
					1                                 & \alpha_{bc}\alpha_{ca} & \alpha_{ca}            & \alpha_{da}                                  & \alpha_{eb}\alpha_{bc}\alpha_{ca} & 0      \\
					\alpha_{ab}                       & 1                      & \alpha_{ca}\alpha_{ab} & \alpha_{da}\alpha_{ab}                       & \alpha_{eb}                       & 0      \\
					\alpha_{ab}\alpha_{bc}            & \alpha_{bc}            & 1                      & \alpha_{da}\alpha_{ab}\alpha_{bc}            & \alpha_{eb}\alpha_{bc}            & 0      \\
					0                                 & 0                      & 0                      & \Delta                                       & 0                                 & 0      \\
					0                                 & 0                      & 0                      & 0                                            & \Delta                            & 0      \\
					\alpha_{ab}\alpha_{bc}\alpha_{cf} & \alpha_{bc}\alpha_{cf} & \alpha_{cf}            & \alpha_{da}\alpha_{ab}\alpha_{bc}\alpha_{cf} & \alpha_{eb}\alpha_{bc}\alpha_{cf} & \Delta \\
				\end{pmatrix}
				\begin{pmatrix}
					\varepsilon_A \\ \varepsilon_B \\ \varepsilon_C \\ \varepsilon_D \\ \varepsilon_E \\ \varepsilon_F
				\end{pmatrix}
			\end{equation*}
			\item \points{1} E.g. intervene $X_A=x_a$. Then in $M$, $B$ is affected by $x_a$, but in $M^{\text{acy}}$, it is not. This should be proven formally.

			In $M_{\doit(X_A = x_A)}$ we have solution function
			\begin{equation*}
				g_B(x_A, \varepsilon) = \alpha_{ab}x_A + \alpha_{eb}\varepsilon_E
			\end{equation*}
			so $P_M(X_B\given \doit(X_A=x_A)) = \Ncal(\alpha_{ab}x_A + \alpha_{eb}\mu_E, \alpha_{eb}^2\sigma_{E}^2)$.

			In $M^{\text{acy}}_{\doit(X_A = x_A)}$ we have solution function
			\begin{equation*}
				\tilde{g}_B(x_A, \varepsilon) = \frac{1}{\Delta}(\alpha_{ab}\varepsilon_A + \varepsilon_B + \alpha_{ca}\alpha_{ab}\varepsilon_C + \alpha_{da}\alpha_{ab}\varepsilon_D + \alpha_{eb}\varepsilon_E)
			\end{equation*}
			so $P_{M^{\text{acy}}}(X_B\given \doit(X_A=x_A)) = \Ncal(\mu, \sigma^2)$ with
			\begin{align*}
				\mu    & = \frac{1}{\Delta}(\alpha_{ab}\mu_A + \mu_B + \alpha_{ca}\alpha_{ab}\mu_C + \alpha_{da}\alpha_{ab}\mu_D + \alpha_{eb}\mu_E) \\
				\sigma & = \frac{1}{\Delta^2}(\alpha_{ab}^2\sigma_A + \sigma_B + \alpha_{ca}^2\alpha_{ab}^2\sigma_C + \alpha_{da}^2\alpha_{ab}^2\sigma_D + \alpha_{eb}^2\sigma_E),
			\end{align*}
			so
			\begin{equation*}
				P_M(X_B\given \doit(X_A=x_A)) \neq P_{M^{\text{acy}}}(X_B\given \doit(X_A=x_A)).
			\end{equation*}
			It suffices to remark that the solution function for $\tilde{g}_B$ in $M^{\text{acy}}$ does not depend on $x_A$ without explicit calculations.
			\item \points{3} The first two are false, the third is true. Prove this!
		\end{enumerate}
	\end{gradedsolution}

\end{enumerate}
\end{document}
